\begin{abstract}
We study what an explicit construction would have to establish for the
Geometric Selection Conjecture. We give a smooth globally proper selection
functional once a compatible viable point is supplied, distinguish this
marked-point construction from a selector intrinsic to an unmarked metric,
and prove a quantitative criterion for simultaneous stabilization and
preservation of physical mass inequalities. We independently check the
slope, Chern-class and index data of a published tetraquadric line-bundle
example. An explicit normalization example shows that a holomorphic Yukawa
hierarchy need not be a physical hierarchy. A local supersymmetric effective
theory has an isolated minimum with broken supersymmetry, but its realization
by the compactification is not established. No vacuum satisfying all nine
conditions is constructed, and the full conjecture remains open.
\end{abstract}
\maketitle

\noindent\textbf{Research status.} This is an author-review draft for a
possible revision of \cite{Mirror}. It is not a proof of the full Geometric
Selection Conjecture (GSC). The use of an Annals manuscript class concerns
formatting; no journal acceptance, submission, or new Hexagon identifier
is asserted. The elementary construction criteria below are not claimed
to be previously unknown theorems.

\section{The simultaneous construction problem}
The GSC formulation in \cite{Dwork,Mirror} concerns a landscape
$\mathcal L=\coprod_\tau M_\tau$ of compactification data and an involutive
isometry $\mu$ defined on a union $D$ of components. A desired function
$S:\mathcal L\to[0,\infty)$ is smooth and proper on each component,
has a unique global minimum at a viable point, and is invariant under
$\mu$ wherever $\mu$ is defined. Its viability requirements originate in
\cite{Original}. We use the following short restatement, without supplying
new numerical meanings for the physical acceptance conditions.

\begin{center}\small
\begin{tabular}{@{}p{.07\linewidth}p{.85\linewidth}@{}}
\toprule
C1 & A slope-polystable bundle with $c_1=0$ and an effective anomaly class.\\
C2 & Three net families, with the quotient index treated consistently.\\
C3 & Wilson-line reduction of the visible gauge group to the Standard Model.\\
C4 & Three quark--lepton families and no chiral exotics after projection.\\
C5 & Allowed sufficiently heavy masses for vectorlike exotics.\\
C6 & A Higgs pair with doublet--triplet splitting.\\
C7 & Stabilization of the relevant moduli in a controlled effective regime.\\
C8 & Acceptable physical Yukawa rank and hierarchy after normalization.\\
C9 & Phenomenologically consistent supersymmetry breaking, if present.\\
\bottomrule
\end{tabular}
\end{center}

These are nine simultaneous predicates on one complete datum, not nine
alternative cases. Let $A_i$ be the locus satisfying C$i$. The existence
problem is $\bigcap_{i=1}^9 A_i\ne\varnothing$. Nonemptiness of separate
$A_i$, even of all pairs among three loci, need not imply a common point:
$\{0,1\}$, $\{1,2\}$ and $\{0,2\}$ have pairwise intersections but empty
triple intersection. Results at different bundles or different moduli
values cannot be assembled into a single vacuum without proving their
compatibility.

For a complete datum one must specify the threefold, visible and hidden
bundles, polarization, free action and equivariant structures when used,
Wilson line, effective K\"ahler potential and superpotential, and the
domain in which omitted corrections are controlled. C5 and C7--C9 also
need explicit acceptance predicates. We leave those predicates as inputs;
choosing convenient thresholds would change the problem.

\section{A marked-point selection construction}
The functional conditions have a clean differential-topological answer.
They do not manufacture the point that is to be selected.

\begin{theorem}[Selection with a supplied point]\label{thm:selection}
Let $M$ be a countable disjoint union of second-countable smooth
Riemannian manifolds. Let $\mu$ be an involutive isometry on a union $D$
of components. Given $p\in M$, the following are equivalent:
\begin{enumerate}[label=(\roman*)]
\item $p\notin D$, or $\mu p=p$.
\item There is a nonnegative smooth globally proper function $S$ on $M$
with unique zero $p$ and $S\circ\mu=S$ on $D$.
\end{enumerate}
The function can be chosen to have Hessian $2g_p$ at $p$.
\end{theorem}
\begin{proof}
An invariant function with unique zero must fix that zero under every
symmetry defined there. This proves necessity.

Every second-countable smooth manifold has a nonnegative smooth proper
exhaustion $\rho$. One direct construction takes a compact exhaustion
$K_j\subset\operatorname{int}K_{j+1}$, a locally finite partition of unity
subordinate to the compact annular cover, and assigns successively increasing
positive coefficients to its terms; outside $K_j$ all nonzero coefficients
can be arranged to be at least $j$. Its sublevel sets are closed and
contained in some $K_j$, hence compact. These standard smooth-manifold
constructions apply separately to the components; see \cite{Lee}.

On a component preserved by $\mu$, replace $\rho$ by
$\rho+\rho\circ\mu$. On a pair of exchanged components choose $\rho$ on
one and transport it to the other. The resulting function is invariant
and proper on each component. Enumerate the component orbits; outside
the orbit of the component of $p$, add distinct positive integer offsets
tending to infinity. Each orbit contains at most two components.

Choose a normal ball $B_r(p)$ with compact closure, and write
$f(x)=d(p,x)^2$ there. If $p\in D$ and is fixed, this ball and $f$ are
invariant. Choose a smooth function $\theta:[0,\infty)\to[0,1]$ equal to
one for $t\le r^2/4$ and zero for $t\ge r^2/2$. On the component of $p$ set
\[
 S=\theta(f)f+(1-\theta(f))(1+\rho),
\]
extending the expression by $1+\rho$ outside the ball. This is smooth
because $\theta(f)$ vanishes near its boundary. Its only zero is $p$:
away from $p$ every nonzero term in this convex combination is positive.
It equals $f$ near $p$, so its Hessian there is $2g_p$.

On the remaining components use their exhaustion plus the positive
orbit offset. For a fixed sublevel only finitely many offsets occur.
On the distinguished component a sublevel is a closed subset of
$\overline B_r(p)\cup\{\rho\le t-1\}$; on all other relevant components
it is compact by properness of $\rho$. Thus every sublevel of $S$ is
compact in the disjoint union. Invariance follows from the construction.
\end{proof}

\begin{remark}[What the point supplies]
The construction uses the marked point $p$, exhaustions and cutoffs in
addition to the metric. It proves the smooth-function assertion with a
supplied compatible point. It does not prove that an unmarked metric
determines $p$ or that $p$ satisfies C1--C9.
\end{remark}

\begin{proposition}[An intrinsic selector must respect all isometries]
Suppose a rule assigns $S_g$ to an unmarked Riemannian manifold naturally
under isometries. If $S_g$ has a unique minimum $p$, every isometry fixes
$p$. In particular there is no such uniquely minimizing rule on a
positive-dimensional homogeneous Riemannian manifold.
\end{proposition}
\begin{proof}
Naturality for an isometry $h$ says $S_g(hx)=S_g(x)$. Hence $hp$ is
another minimum and equals $p$. Transitivity on a manifold with more
than one point is incompatible with a common fixed point.
\end{proof}
This is an obstruction to a particular notion of metric-only selection,
not a proof that the actual heterotic moduli metric is homogeneous.
It also shows why invariance under an additional compact group cannot be
discarded merely because the partial mirror map is undefined at $p$:
the additional group must still fix a unique selected point.

\section{Index constraints and the mirror restriction}
\begin{proposition}[Index and finite descent]\label{prop:index}
Let $X$ be a Calabi--Yau threefold and $V$ a holomorphic bundle with
$c_1(V)=0$. Then
\[
 \chi(X,V)=\frac12\int_Xc_3(V).
\]
If $\pi:X\to Y$ is finite \emph{\'etale} of degree $d$ and
$V=\pi^*V_Y$, then $\chi(X,V)=d\chi(Y,V_Y)$. Thus the three-family
index condition downstairs requires $|\int_Xc_3(V)|=6d$.
\end{proposition}
\begin{proof}
In Hirzebruch--Riemann--Roch, the degree-six term of
$\operatorname{ch}(V)\operatorname{td}(X)$ is
$\operatorname{ch}_3(V)+c_1(V)c_2(X)/12+operatorname{rk}(V)c_1(X)c_2(X)/24$.
The last two terms vanish and $\operatorname{ch}_3(V)=c_3(V)/2$.
For an \emph{\'etale} map the tangent bundle and Todd class pull back,
and integration of a pulled-back top class multiplies by $d$.
\end{proof}

For a quintic, expansion of $(1+H)^5/(1+5H)$ gives
$c_3(T_X)=-40H^3$ and $\int_XH^3=5$, hence $\int c_3(T_X)=-200$.
Consequently no standard-embedding
free quotient has index of absolute value three, because $100=3d$ has no
integer solution. This restriction is independent of the quotient group
classification. It does not apply to arbitrary nonstandard bundles.

Under the standard-type partial mirror convention of \cite{Mirror}, a
fixed point has $h^{1,1}=h^{2,1}$ and hence $\chi(X)=0$. The bundle on
that component has $\int c_3(V)=\chi(X)$, so Proposition~\ref{prop:index}
gives zero net families. This proves that a selected viable point cannot
be a mirror-fixed point in that specified domain. It does not exclude
nonstandard data outside the domain, nor establish a mirror map on all
$(0,2)$ compactifications.

\section{An exact nonstandard geometric starting point}
We now check a published bundle example rather than assume that a
topological index determines its masses. The bundle data are those of
Buchbinder, Constantin and Lukas \cite[Sec.~5.1]{BCL}; the calculations
below are independent reproductions, not a new discovery of their model.

Let $X\subset A=(\mathbb P^1)^4$ be a smooth hypersurface of multidegree
$(2,2,2,2)$. Write $J_i$ for the four hyperplane classes. Adjunction gives
$K_X\cong\mathcal O_X$ and
\[
 d_{ijk}=\int_XJ_iJ_jJ_k=
 \begin{cases}2& i,j,k\text{ distinct},\\0&\text{otherwise},\end{cases}
 \qquad c_2(T_X)=4\sum_{i<j}J_iJ_j.
\]
These formulas follow in the ambient ring $J_i^2=0$ from
$[X]=2\sum J_i$ and $c(T_X)=\prod_i(1+2J_i)/(1+2\sum J_i)$.
Set $V=\bigoplus_{a=1}^5\mathcal O_X(k_a)$, with columns
\begin{equation}\label{eq:bundle}
(k_a^i)=
\begin{pmatrix}
-1&-1&0&1&1\\
0&-3&1&1&1\\
0&2&-1&-1&0\\
1&2&0&-1&-2
\end{pmatrix}.
\end{equation}

\begin{proposition}[Cover-space data]\label{prop:cover}
For every smooth $X$ above, $c_1(V)=0$ and $V$ is slope-polystable for
$J=t(J_1+J_2+J_3+J_4)$, $t>0$. Its Chern and index data are
\[
 \left(\int_Xc_2(V)J_i\right)_{i=1}^4=(24,8,20,12),
 \qquad \int_Xc_3(V)=-24,\qquad\chi(X,V)=-12.
\]
The anomaly class on the cover has the effective representative
$6J_1J_3+2J_1J_4$.
\end{proposition}
\begin{proof}
Every row of \eqref{eq:bundle} sums to zero, proving $c_1(V)=0$.
Every column also sums to zero. The slope of a summand is therefore
$12t^2\sum_i k_a^i=0$. Rank-one bundles are slope-stable, so their
direct sum at common slope zero is polystable.

Put $l_a=\sum_i k_a^iJ_i$. From $\sum l_a=0$ one obtains
$c_2(V)=-\tfrac12\sum l_a^2$ and
$c_3(V)=\tfrac13\sum l_a^3$. Direct multiplication using $d_{ijk}$ gives
the displayed values. Independently, for a line bundle,
\[
 \chi(\mathcal O_X(k))=
 2\sum_{i<j<k'}k_i k_j k_{k'}+2\sum_i k_i.
\]
For the five columns the indices are $0,-8,0,0,-4$, giving $-12$.
The anomaly pairs with the $J_i$ as $(0,16,4,12)$, exactly the pairings
of $6J_1J_3+2J_1J_4$. The $J_i$ form an integral basis of $H^2(X)$ by
Lefschetz; $H^4(X)$ is torsion-free by Lefschetz on $H_2$ and Poincar\'e
duality. Equality of these pairings therefore identifies the classes.
Each $J_iJ_j$ with $i\ne j$ is represented by an effective intersection
of divisors, proving the assertion.
\end{proof}

\subsection{A free degree-four family and equivariance}
On every factor let $a[x:y]=[x:-y]$ and $b[x:y]=[y:x]$ and act
diagonally on $A$. They commute projectively and generate
$\Gamma\cong(\mathbb Z/2)^2$. Each nonidentity element fixes $2^4=16$
ambient points. The three sets are disjoint, so there are 48 points to
avoid. Put
\[
 e_i=x_i^2+y_i^2,\quad d_i=x_i^2-y_i^2,\quad p_i=2x_iy_i.
\]
An invariant multidegree-$(2,2,2,2)$ linear system is spanned by products
of one of $e_i,d_i,p_i$ per factor with even numbers of $d$ and $p$.

\begin{lemma}[Base-point-free invariant system]
This invariant linear system is base-point-free. Its general member is
smooth and avoids all 48 fixed points, so $\Gamma$ acts freely on it.
\end{lemma}
\begin{proof}
At a point of any factor at least two of $e,d,p$ are nonzero. Label them
by $0,u,v\in(\mathbb Z/2)^2$. Any two-element subset of these three labels
is an affine line. If the available two-element sets in four factors
have different directions, their sumset is the whole group. If all have
the same direction, they are the same pair and their fourfold sumset
contains zero. A choice of nonzero factors with total label zero therefore
exists, proving base-point-freeness. Taking subsets deals also with a
factor where all three values are nonzero.

Bertini gives smoothness on a nonempty open subset of the parameter
space. At any of the finitely many fixed points base-point-freeness
makes the vanishing condition a proper hyperplane. The complement of
these hyperplanes meets that open subset. Such a member has a free
action. The product of the four determinants of either ambient lift is
$(-1)^4=1$, and its defining polynomial is invariant, so the residue
holomorphic three-form is preserved. Thus the free quotient has trivial
canonical bundle as well.
\end{proof}

The two linear lifts anticommute on $\mathcal O(1)$ of one factor.
Their commutator on $\mathcal O(k_1,k_2,k_3,k_4)$ is
$(-1)^{\sum k_i}$. Every column of \eqref{eq:bundle} has sum zero, so
these tensor lifts commute and square to one. They give equivariant
structures on the summands and a descended bundle $V_Y$ on $Y=X/\Gamma$.
Proposition~\ref{prop:index} now proves $\chi(Y,V_Y)=-3$.

The invariant polarization descends and proves polystability downstairs.
The anomaly class is rationally effective downstairs: pushing forward
an effective representative upstairs and dividing by four proves this
over $\mathbb Q$. An equality after pullback alone does not exclude
integral torsion in $H^4(Y,\mathbb Z)$. We have not audited that torsion
refinement or the hidden-sector completion here. We therefore distinguish
the verified rational anomaly statement from a full quantized anomaly
certificate for the quotient.

A Wilson homomorphism on the $SU(5)$ factor may send $a$ to
$\operatorname{diag}(1,1,1,-1,-1)$ and $b$ to the identity. Its centralizer
is $S(U(3)\times U(2))$, the Standard Model group up to a finite quotient.
This group calculation is not a calculation of all projected cohomology
groups, possible additional gauge factors, or their masses. Those are
separate C3--C6 obligations. The example is a geometric starting point,
not a certificate of a nine-condition vacuum.

\section{Why holomorphic masses do not close C8}
After fixing conventions and the Higgs factor, a physical Yukawa matrix
has the form
\begin{equation}\label{eq:normalisation}
 Y_{\rm phys}=c\,Z_L^{-T/2}Y_{\rm hol}Z_R^{-1/2},
 \qquad c>0,
\end{equation}
where the matter metrics are positive definite. Their invertibility
preserves rank but does not preserve singular-value ratios.

\begin{proposition}[A normalization counterexample]\label{prop:normalisation}
For $0<\epsilon<1$, take
\[
Y_{\rm hol}=\operatorname{diag}(\epsilon^8,\epsilon^4,1),\quad
Z_L=I,\quad Z_R=\operatorname{diag}(\epsilon^{16},\epsilon^8,1),\quad c=1.
\]
All metrics are positive definite and $Y_{\rm phys}=I$. Thus a full-rank
holomorphic matrix with separated scales does not prove a physical hierarchy.
\end{proposition}
\begin{proof}
The positive inverse square root of $Z_R$ is
$\operatorname{diag}(\epsilon^{-8},\epsilon^{-4},1)$. Multiplication in
\eqref{eq:normalisation} gives the identity.
\end{proof}
The example does not claim that these metrics occur in the tetraquadric
model. It refutes the inference based on a holomorphic matrix alone.
It does not refute a hierarchy proved with regular nondegenerate metrics
at the limiting point, as the following bound makes precise.

\begin{proposition}[Bounded normalization preserves orders]\label{prop:bounded}
Suppose $mI\preceq Z_L,Z_R\preceq MI$ with $0<m\le M$ and
$0<c_-\le c\le c_+$. For every singular value, ordered consistently,
\[
 \frac{c_-}{M}\sigma_i(Y_{\rm hol})
 \le\sigma_i(Y_{\rm phys})
 \le\frac{c_+}{m}\sigma_i(Y_{\rm hol}).
\]
If these bounds are uniform as a parameter tends to zero, any nonzero
power-law orders of the holomorphic singular values are preserved.
\end{proposition}
\begin{proof}
For invertible $A,B$, the min--max characterization gives
$\sigma_i(AYB)\le\|A\|_{\rm op}\|B\|_{\rm op}\sigma_i(Y)$.
Apply this also to $Y=A^{-1}(AYB)B^{-1}$ to obtain the reverse bound.
The operator norms of the inverse square roots are at most $m^{-1/2}$
and those of their inverses at most $M^{1/2}$; transpose has the same
singular values. Multiplication by $c$ yields the result. Uniform positive
constants do not change a power-law order.
\end{proof}
Continuity and positive definiteness at a regular limiting point give
such bounds locally. The preceding counterexample loses that lower bound
as $\epsilon\to0$. Conversely, an asymptotic order calculation alone does
not fix the observed mass ratios or the point selected by stabilization.
Published physical-flavour constructions are available \cite{CLLN};
that paper leaves stabilization and supersymmetry breaking untreated.
This is a limitation of the result cited, not an assertion about every
heterotic model in the literature.

\section{A quantitative simultaneous-persistence criterion}
The following theorem identifies bounds that would make a local
construction robust. All its assumptions must be checked on the same
compactification chart and in the same normalization.

\begin{theorem}[A certified interior minimum]\label{thm:minimum}
Let $B=\overline B_r(u_0)\subset\mathbb R^n$ and let $V_0,\Delta$
be $C^2$ on a neighborhood of $B$. Suppose on $B$ that
\[
 D^2V_0\succeq\lambda I,\quad
 \|D^2\Delta\|_{\rm op}\le b<\lambda,\quad
 \|\nabla V_0(u_0)\|\le\eta,\quad
 \|\nabla\Delta(u_0)\|\le a.
\]
Put $\kappa=\lambda-b$ and assume $\eta+a<\kappa r$.
Then $V=V_0+\Delta$ has exactly one critical point $u_*$ in $B$.
It is an interior minimum, $D^2V(u_*)\succeq\kappa I$, and
\begin{equation}\label{eq:displacement}
 \|u_*-u_0\|\le\delta:=\frac{\eta+a}{\kappa}<r.
\end{equation}
\end{theorem}
\begin{proof}
The Hessian of $V$ is bounded below by $\kappa I$. Integrating along
a segment yields strong monotonicity:
\[
 (\nabla V(x)-\nabla V(y))\cdot(x-y)\ge\kappa\|x-y\|^2.
\]
For $x$ on the boundary and $v=x-u_0$ this gives
$\nabla V(x)\cdot v\ge\kappa r^2-(\eta+a)r>0$.
The inward directional derivative is negative, so a minimum on the
compact ball cannot occur at its boundary. A minimum exists and is
therefore critical in the interior. Strong monotonicity excludes two
critical points. Apply it with $x=u_*$, $y=u_0$ and use
$\nabla V(u_*)=0$ to obtain
$\kappa\|u_*-u_0\|^2\le(\eta+a)\|u_*-u_0\|$,
which proves \eqref{eq:displacement}.
\end{proof}

\begin{proposition}[Preserving masses and inequalities]\label{prop:persistence}
In Theorem~\ref{thm:minimum}, suppose $Y$ is a physical $3\times3$ Yukawa
matrix with $\|Y(u)-Y(u_0)\|_{\rm op}\le L\|u-u_0\|$ on $B$.
If $s_i$ are its singular values at $u_0$, then at $u_*$ they differ
from $s_i$ by at most $e=L\delta$. In particular $s_{\min}>e$
certifies full rank. For $s_j>e$,
\[
 \frac{\max(0,s_i-e)}{s_j+e}
 \le\frac{s_i(u_*)}{s_j(u_*)}
 \le\frac{s_i+e}{s_j-e}.
\]
If a real acceptance function $q$ is $L_q$-Lipschitz and has required
lower bound $q_{\rm req}$, then
$q(u_0)>q_{\rm req}+L_q\delta$ certifies $q(u_*)>q_{\rm req}$.
\end{proposition}
\begin{proof}
The min--max formula for singular values and
$|\|Av\|-\|Bv\||\le\|(A-B)v\|$ on unit vectors give
$|\sigma_i(A)-\sigma_i(B)|\le\|A-B\|_{\rm op}$ by applying the
formula in both directions. The ratio bounds follow by bounding numerator
and denominator separately. The last assertion is the Lipschitz inequality.
\end{proof}

Thus a controlled ball preserving the discrete spectrum and satisfying
strict mass, hierarchy and consistency margins would turn an approximate
minimum into a simultaneous certificate. If the real kinetic matrix obeys
$0<G(u_*)\preceq\Lambda I$, the physical scalar mass matrix
$G(u_*)^{-1/2}D^2V(u_*)G(u_*)^{-1/2}$ has eigenvalues at least
$\kappa/\Lambda$. This follows by applying the Hessian lower bound to
$G^{-1/2}v$. At a critical point the connection correction to the Hessian
vanishes. Thus kinetic normalization must also be bounded when interpreting
the Hessian physically. The bounds must include all
claimed perturbative and nonperturbative corrections, the normalization
metrics, the physical moduli coordinates, and any constraints or gauge
quotients. A Hessian estimate for a freely invented scalar function is
not a Hessian estimate for the compactification potential.

\section{An explicit local effective-theory construction and its limit}
It is useful to separate a solvable analytic problem from the missing
string realization. Consider a classical global $N=1$ effective theory
with fields $S,Z_1,\ldots,Z_m$, constants $a>0$, $F\ne0$, and $m_j\ne0$:
\[
 K=|S|^2-a|S|^4+\sum_j|Z_j|^2,
 \qquad W=FS+\frac12\sum_jm_jZ_j^2.
\]
We work on $|S|^2<1/(4a)$, where its K\"ahler metric is positive.

\begin{proposition}[A stabilized local model]\label{prop:eft}
For this specified classical global effective theory,
\[
 V=\frac{|F|^2}{1-4a|S|^2}+\sum_j|m_j|^2|Z_j|^2
\]
has a unique minimum at the origin of the stated domain. Supersymmetry
is broken there. In canonical real coordinates its Hessian has diagonal
entries $4a|F|^2,4a|F|^2$ and $|m_j|^2,|m_j|^2$ for each $j$.
\end{proposition}
\begin{proof}
$K_{S\bar S}=1-4a|S|^2$ and $K_{Z_i\bar Z_j}=\delta_{ij}$, so the
global-supersymmetry potential $K^{A\bar B}W_A\overline W_B$ is the
displayed expression. Its first term is strictly increasing with $|S|^2$
and the remaining terms are positive quadratics. The auxiliary field in
the $S$ direction is nonzero at zero because $W_S=F$. Expanding with
$S=(x+iy)/\sqrt2$ gives
$V=|F|^2+2a|F|^2(x^2+y^2)+\sum_j|m_j|^2(u_j^2+v_j^2)/2+O(|S|^4)$,
which proves the Hessian assertion.
\end{proof}
This proposition supplies no Calabi--Yau bundle realizing its $K$ and $W$.
It omits supergravity, flux quantization, the anomaly constraints and
the coupled visible sector. Consequently it verifies neither C7 nor C9
for the geometric example. Its role is to make the missing realization
problem explicit rather than hide it inside a choice of coefficients.

\begin{proposition}[A tree-level obstruction]\label{prop:noscale}
Let a supergravity model have
\[
 K=K_0(z,\bar z)-\sum_{i=1}^r n_i\log(T_i+\bar T_i),
 \quad n_i>0,\quad\sum_i n_i=3,\quad W=W(z).
\]
At a point with $D_zW=0$, its F-term potential vanishes for all $T_i$
in this product chart. It does not stabilize these K\"ahler directions.
\end{proposition}
\begin{proof}
The $T$-block is diagonal,
$K_i=-n_i/(T_i+\bar T_i)$ and
$K_{i\bar i}=n_i/(T_i+\bar T_i)^2$. Hence
$\sum_iK^{i\bar i}K_iK_{\bar i}=\sum_i n_i=3$.
Since $D_iW=K_iW$ and the $z$ block contributes zero, these terms cancel
the $-3|W|^2$ in $V_F=e^K(K^{A\bar B}D_AW\overline{D_BW}-3|W|^2)$.
Thus there are flat $T$ directions. Additional D terms or corrections are
not covered by this assertion and would need their own calculation.
\end{proof}

\section{The remaining proof obligations}
The geometric calculation supplies explicit nonstandard bundle data,
cover-space polystability and anomaly effectiveness, and an exact
three-family quotient index. It does not compute the full projected
spectrum or all mass operators. In particular the quotient integral
anomaly certificate and the hidden sector need a complete audit.

The local effective model proves that a specified mathematical potential
can have a positive Hessian and broken supersymmetry. It does not show that
this potential and the required physical Yukawa matrices arise from the
same compactification. Theorem~\ref{thm:minimum} provides a route for a
controlled candidate: establish its hypotheses and the acceptance margins
on that candidate's chart. Those estimates have not been supplied here.

Even an explicit metric-dependent selection function, a correct index,
or a holomorphic hierarchy cannot replace these missing results. No
intersection point of all nine viability loci has been established by
the present constructions. The full GSC remains open in this draft.

\section*{Verification and provenance}
Exact-arithmetic controls reproduce the Chern, slope and index data, the
256 invariant-system support cases, and the normalization and minimum
examples. Sympy checks the local Hessian and no-scale identity. These are
algebra checks, not a nine-condition construction; no Lean formalization
is claimed.

ChatGPT/Codex assisted with this draft at the author's request. Both
existing authors are retained, and their review of the new material is
pending. No v2 has been uploaded to Hexagon.

\begin{thebibliography}{99}
\bibitem{Original} D.~Bhattacharjee, U.~Bhattacharya and S.~Ghosh,
\emph{One missing axiom for the complete closure of the heterotic landscape},
Preprints.org (2026), version 2,
\href{https://doi.org/10.20944/preprints202603.0792.v2}{doi:10.20944/preprints202603.0792.v2}.
Cited for the formulation of the nine conditions, not as a proved
existence theorem.
\bibitem{Dwork} D.~Bhattacharjee and U.~Bhattacharya,
\emph{Which string vibrations become which generations on the Dwork quintic},
\href{https://hexagonmath.org/2610.00446v1}{Hexagon 2610.00446v1};
code archive \href{https://github.com/creelie/strings/releases/tag/v1.2.0}{creelie/strings v1.2.0}.
\bibitem{Mirror} D.~Bhattacharjee and U.~Bhattacharya,
\emph{The mirror clause of the Geometric Selection Conjecture is empty},
\href{https://hexagonmath.org/2610.00493v1}{Hexagon 2610.00493v1};
\href{https://doi.org/10.5281/zenodo.23567115}{doi:10.5281/zenodo.23567115}.
\bibitem{BCL} E.~I.~Buchbinder, A.~Constantin and A.~Lukas,
\emph{The moduli space of heterotic line bundle models: a case study for the
tetra-quadric}, J. High Energy Phys. \textbf{03} (2014), 025,
\href{https://doi.org/10.1007/JHEP03(2014)025}{doi:10.1007/JHEP03(2014)025}.
The bundle in \eqref{eq:bundle} is the published example of Section 5.1.
\bibitem{CLLN} A.~Constantin, L.~T.-Y.~Leung, A.~Lukas and L.~A.~Nutricati,
\emph{Reproducing Standard Model fermion masses and mixing in string theory:
a heterotic line bundle study}, Phys. Rev. D \textbf{113} (2026), 046005,
\href{https://doi.org/10.1103/fbdz-t73z}{doi:10.1103/fbdz-t73z}.
\bibitem{Hartshorne} R.~Hartshorne, \emph{Algebraic Geometry},
Graduate Texts in Mathematics 52, Springer, 1977.
\bibitem{Lee} J.~M.~Lee, \emph{Introduction to Smooth Manifolds},
second ed., Graduate Texts in Mathematics 218, Springer, 2013.
\end{thebibliography}
